
\section{La célula electroquímica y la geometría del ataque}
Una célula de corrosión necesita procesos anódico y catódico, conexión electrónica y un medio que permita transporte iónico. En una pareja galvánica importan potenciales en el ambiente real, polarización, área y resistencia del circuito. Un ánodo pequeño conectado a un cátodo grande puede sufrir una densidad de corriente anódica elevada. Invertir las áreas cambia la severidad aunque permanezca la misma pareja nominal.

Los huecos retienen electrolito y favorecen diferencias de aireación; un ataque localizado puede progresar sin una pérdida de masa total llamativa. Por ello la inspección de uniones, zonas de drenaje y superficies ocultas se relaciona con geometría y acceso. Un acabado protege mientras conserva continuidad y compatibilidad; una discontinuidad puede concentrar ataque en una zona pequeña.
\section{Daño ambiental que interactúa con tensiones}
La corrosión bajo tensión requiere una combinación susceptible de material y estado, ambiente y tensión de tracción, incluida la residual. No es simplemente corrosión general acelerada por cualquier carga. La fatiga--corrosión combina ciclos y ambiente y puede modificar iniciación, umbral y crecimiento; las curvas en aire no se trasladan sin más a ese entorno. En ambos casos es esencial identificar mecanismo y condiciones, no solo medir una dimensión final del defecto.

La fragilización por hidrógeno puede relacionarse con procesos de fabricación, protección o servicio, según material y exposición. Un tratamiento o protección que mejore una característica no garantiza ausencia de efectos secundarios. El análisis de fallo examina procedencia, microestructura, superficie de fractura y secuencia temporal; una sola fotografía raramente prueba la causa.
\section{Oxidación protectora y pérdida de protección}
Una película compacta y adherente puede frenar difusión y dar una cinética aproximadamente parabólica. Una película porosa, fisurada o que se desprende puede no mantener esa protección. Los ciclos térmicos introducen incompatibilidades de dilatación y tensiones; el espesor de óxido no equivale directamente a espesor de metal consumido sin conocer composición, densidades y geometría.

En zonas calientes también puede haber corrosión por depósitos y especies ambientales que alteran la película. No se agrupan todos los ataques de alta temperatura bajo un único parámetro de oxidación. En mantenimiento, limpieza e inspección deben ajustarse a datos del producto: una retirada agresiva de depósito puede dañar el sustrato o borrar evidencia.
\repaso{¿Una unión carbono--aluminio es segura frente a corrosión por no contener dos metales?}{No puede concluirse así. El carbono conductor puede participar en un acoplamiento electroquímico con el aluminio en un ambiente adecuado. Se estudian aislamiento, protección, sellado, drenaje y mantenimiento dentro del diseño aplicable.}
